% 题证摘录；来源与整理边界见相同编号分题文件。

\begin{lemma}[37.53.1]
Let $S$ be a scheme. Let $f:X\to S$ be a universally closed and quasi-separated morphism. There exists a factorization
\[\xymatrix{X\ar[rr]_{f'}\ar[rd]_f&&S'\ar[dl]^\pi\\&S&}\]
with the following properties:
\begin{enumerate}
\item the morphism $f'$ is universally closed, quasi-compact, quasi-separated, and surjective,
\item the morphism $\pi:S'\to S$ is integral,
\item we have $f'_*\mathcal O_X=\mathcal O_{S'}$,
\item we have $S'=\underline{\Spec}_S(f_*\mathcal O_X)$, and
\item $S'$ is the normalization of $S$ in $X$, see Morphisms, Definition 29.54.3.
\end{enumerate}
Formation of the factorization $f=\pi\circ f'$ commutes with flat base change.
\end{lemma}
\begin{proof}
By Morphisms, Lemma 29.42.8 the morphism $f$ is quasi-compact. Hence the normalization $S'$ of $S$ in $X$ is defined (Morphisms, Definition 29.54.3) and we have the factorization $X\to S'\to S$. By Morphisms, Lemma 29.54.11 we have (2), (4), and (5). The morphism $f'$ is universally closed by Morphisms, Lemma 29.42.7. It is quasi-compact by Schemes, Lemma 26.21.14 and quasi-separated by Schemes, Lemma 26.21.13.

To show the remaining statements we may assume the base scheme $S$ is affine, say $S=\Spec(R)$. Then $S'=\Spec(A)$ with $A=\Gamma(X,\mathcal O_X)$ an integral $R$-algebra. Thus it is clear that $f'_*\mathcal O_X$ is $\mathcal O_{S'}$ (because $f'_*\mathcal O_X$ is quasi-coherent, by Schemes, Lemma 26.24.1, and hence equal to $\widetilde A$). This proves (3).

Let us show that $f'$ is surjective. As $f'$ is universally closed (see above) the image of $f'$ is a closed subset $V(I)\subset S'=\Spec(A)$. Pick $h\in I$. Then $h|_X=f^\sharp(h)$ is a global section of the structure sheaf of $X$ which vanishes at every point. As $X$ is quasi-compact this means that $h|_X$ is a nilpotent section, i.e., $h^n|X=0$ for some $n>0$. But $A=\Gamma(X,\mathcal O_X)$, hence $h^n=0$. As every element of $I$ is nilpotent, we conclude that $V(I)=S'$ as desired.

By Cohomology of Schemes, Lemma 30.5.2 we see that formation of $f_*\mathcal O_X$ commutes with flat base change. Formation of the relative spectrum commutes with any base change by Constructions, Lemma 27.4.6. Thus formation of the factorization commutes with flat base change.
\end{proof}

\begin{lemma}[37.53.3]
Let $f:X\to S$ be a morphism of schemes. Let $s\in S$. Then $X_s$ is geometrically connected, if and only if for every \'etale neighbourhood $(U,u)\to(S,s)$ the base change $X_U\to U$ has connected fibre $X_u$.
\end{lemma}
\begin{proof}
If $X_s$ is geometrically connected, then any base change of it is connected. On the other hand, suppose that $X_s$ is not geometrically connected. Then by Varieties, Lemma 33.7.11 we see that $X_s\times_{\Spec(\kappa(s))}\Spec(k)$ is disconnected for some finite separable field extension $k/\kappa(s)$. By Lemma 37.35.2 there exists an affine \'etale neighbourhood $(U,u)\to(S,s)$ such that $\kappa(u)/\kappa(s)$ is identified with $k/\kappa(s)$. In this case $X_u$ is disconnected.
\end{proof}

\begin{theorem}[37.53.4; Stein factorization, Noetherian case]
Let $S$ be a locally Noetherian scheme. Let $f:X\to S$ be a proper morphism. There exists a factorization
\[\xymatrix{X\ar[rr]_{f'}\ar[rd]_f&&S'\ar[dl]^\pi\\&S&}\]
with the following properties:
\begin{enumerate}
\item the morphism $f'$ is proper with geometrically connected fibres,
\item the morphism $\pi:S'\to S$ is finite,
\item we have $f'_*\mathcal O_X=\mathcal O_{S'}$,
\item we have $S'=\underline{\Spec}_S(f_*\mathcal O_X)$, and
\item $S'$ is the normalization of $S$ in $X$, see Morphisms, Definition 29.54.3.
\end{enumerate}
\end{theorem}
\begin{proof}
Let $f=\pi\circ f'$ be the factorization of Lemma 37.53.1. Note that besides the conclusions of Lemma 37.53.1 we also have that $f'$ is separated (Schemes, Lemma 26.21.13) and finite type (Morphisms, Lemma 29.15.8). Hence $f'$ is proper. By Cohomology of Schemes, Proposition 30.19.1 we see that $f_*\mathcal O_X$ is a coherent $\mathcal O_S$-module. Hence we see that $\pi$ is finite, i.e., (2) holds.

This proves all but the most interesting assertion, namely that all the fibres of $f'$ are geometrically connected. It is clear from the discussion above that we may replace $S$ by $S'$, and we may therefore assume that $S$ is Noetherian, affine, $f:X\to S$ is proper, and $f_*\mathcal O_X=\mathcal O_S$. Let $s\in S$ be a point of $S$. We have to show that $X_s$ is geometrically connected.

By Lemma 37.53.3 we see that it suffices to show $X_u$ is connected for every \'etale neighbourhood $(U,u)\to(S,s)$. We may assume $U$ is affine. Thus $U$ is Noetherian (Morphisms, Lemma 29.15.6), the base change $f_U:X_U\to U$ is proper (Morphisms, Lemma 29.42.5), and that also $(f_U)_*\mathcal O_{X_U}=\mathcal O_U$ (Cohomology of Schemes, Lemma 30.5.2).

Hence after replacing $(f:X\to S,s)$ by the base change $(f_U:X_U\to U,u)$ it suffices to prove that the fibre $X_s$ is connected when $f_*\mathcal O_X=\mathcal O_S$. We can deduce this from Derived Categories of Schemes, Lemma 36.32.7 (by looking at idempotents in the structure sheaf of $X_s$) but we will also give a direct argument below.

Namely, we apply the theorem on formal functions, more precisely Cohomology of Schemes, Lemma 30.20.7. It tells us that
\[\mathcal O_{S,s}^{\wedge}=(f_*\mathcal O_X)_s^{\wedge}=\varprojlim_n H^0(X_n,\mathcal O_{X_n})\]
where $X_n$ is the $n$th infinitesimal neighbourhood of $X_s$. Since the underlying topological space of $X_n$ is equal to that of $X_s$ we see that if $X_s=T_1\amalg T_2$ is a disjoint union of nonempty open and closed subschemes, then similarly $X_n=T_{1,n}\amalg T_{2,n}$ for all $n$. And this in turn means $H^0(X_n,\mathcal O_{X_n})$ contains a nontrivial idempotent $e_{1,n}$, namely the function which is identically $1$ on $T_{1,n}$ and identically $0$ on $T_{2,n}$.

It is clear that $e_{1,n+1}$ restricts to $e_{1,n}$ on $X_n$. Hence $e_1=\varprojlim e_{1,n}$ is a nontrivial idempotent of the limit. This contradicts the fact that $\mathcal O_{S,s}^{\wedge}$ is a local ring. Thus the assumption was wrong, i.e., $X_s$ is connected, and we win.
\end{proof}
